% GENERATED FILE. Edit canonical lessons, then run npm run generate:books.
% canonical-source-hash: fnv1a64:76391fe5f70e420c
% canonical-lessons: MR-C478-kavti, MR-C478-kana, MR-C478-bargadi, MR-C478-atde, MR-C478-sandha

\chapter{Bones and Joints}
\label{ch:mr-bones-and-joints}
\hlchaptermodality{478}

\begin{chapteropening}
\textbf{By the end of this chapter:} \emph{I can name the skull, the spine, a rib, the gut and a joint.}

Complete the last lesson of chapter 478: I can name the skull, the spine, a rib, the gut and a joint.
\end{chapteropening}

\section[\texorpdfstring{\mr{कवटी} (kavṭī)}{कवटी (kavṭī)}]{\mr{कवटी} (kavṭī) — the skull}

\label{lesson:MR-C478-kavti}



\begin{quote}
Before the new one: say the Marathi for the calf, then the Marathi for the ankle.
\end{quote}

\subsection*{You'll want to know: \mr{कवटी}}
\textbf{\mr{कवटी}} — \emph{kavṭī} — \textquotedblleft{}the skull\textquotedblright{}.

\begin{quote}
\mr{संग्रहालयात} \mr{एका} \mr{प्राचीन} \mr{माणसाची} \mr{कवटी} \mr{होती}.
\end{quote}

\begin{quote}
\emph{saṅgrahālayāt ekā prācīn māṇsācī kavṭī hotī.}
\end{quote}

\begin{quote}
There was the skull of an ancient man in the museum.
\end{quote}

Feminine.

Feminine; also the shell of a coconut.

\begin{cousinweb}
\textbf{\mr{कवटी}} (\emph{kavṭī}) is a Marathi noun whose origin is not settled.
\end{cousinweb}

\begin{grammarlens}[title={What you've built}]
One more word for describing what you lived and felt.
\end{grammarlens}

\subsection*{Your turn}
\begin{itemize}
\raggedright
  \item \emph{Say it:} \emph{kavṭī}
  \item \emph{Say it:} \emph{kavṭī}, once more
  \item \emph{Say it:} \emph{saṅgrahālayāt ekā prācīn māṇsācī kavṭī hotī.} — There was the skull of an ancient man in the museum.
  \item \emph{Recall:} say the Marathi for the calf, then the Marathi for the ankle, then say \emph{kavṭī} again
\end{itemize}

\begin{tcolorbox}[breakable,colback=teal!4,colframe=teal!35!black,title={Before you move on}]
What does \mr{कवटी} mean? (\textquotedblleft{}The skull\textquotedblright{}.) Say it once more.
\end{tcolorbox}
\section[\texorpdfstring{\mr{कणा} (kaṇā)}{कणा (kaṇā)}]{\mr{कणा} (kaṇā) — the spine; backbone}

\label{lesson:MR-C478-kana}



\begin{quote}
Before the new one: say the Marathi for the ankle, then the Marathi for the skull.
\end{quote}

\subsection*{You'll want to know: \mr{कणा}}
\textbf{\mr{कणा}} — \emph{kaṇā} — \textquotedblleft{}the spine; backbone\textquotedblright{}.

\begin{quote}
\mr{पाठीचा} \mr{कणा} \mr{सरळ} \mr{ठेवून} \mr{तो} \mr{बसला}.
\end{quote}

\begin{quote}
\emph{pāṭhīcā kaṇā saraḷ ṭhevūn to baslā.}
\end{quote}

\begin{quote}
He sat with his spine straight.
\end{quote}

Masculine.

\mr{पाठीचा} \mr{कणा}, the spine; also, figuratively, the backbone of something.

\begin{cousinweb}
\textbf{\mr{कणा}} (\emph{kaṇā}) is a Marathi noun whose origin is not settled.
\end{cousinweb}

\begin{grammarlens}[title={What you've built}]
One more word for describing what you lived and felt.
\end{grammarlens}

\subsection*{Your turn}
\begin{itemize}
\raggedright
  \item \emph{Say it:} \emph{kaṇā}
  \item \emph{Say it:} \emph{kaṇā}, once more
  \item \emph{Say it:} \emph{pāṭhīcā kaṇā saraḷ ṭhevūn to baslā.} — He sat with his spine straight.
  \item \emph{Recall:} say the Marathi for the ankle, then the Marathi for the skull, then say \emph{kaṇā} again
\end{itemize}

\begin{tcolorbox}[breakable,colback=teal!4,colframe=teal!35!black,title={Before you move on}]
What does \mr{कणा} mean? (\textquotedblleft{}The spine; backbone\textquotedblright{}.) Say it once more.
\end{tcolorbox}
\section[\texorpdfstring{\mr{बरगडी} (bargaḍī)}{बरगडी (bargaḍī)}]{\mr{बरगडी} (bargaḍī) — a rib}

\label{lesson:MR-C478-bargadi}



\begin{quote}
Before the new one: say the Marathi for the skull, then the Marathi for the spine.
\end{quote}

\subsection*{You'll want to know: \mr{बरगडी}}
\textbf{\mr{बरगडी}} — \emph{bargaḍī} — \textquotedblleft{}a rib\textquotedblright{}.

\begin{quote}
\mr{पडल्यामुळे} \mr{त्याची} \mr{एक} \mr{बरगडी} \mr{मोडली}.
\end{quote}

\begin{quote}
\emph{paḍlyāmuḷe tyācī ek bargaḍī moḍlī.}
\end{quote}

\begin{quote}
He broke a rib when he fell.
\end{quote}

Feminine. Plural: \mr{बरगड्या}, \emph{bargaḍyā}.

Plural \mr{बरगड्या}, the ribs.

\begin{cousinweb}
\textbf{\mr{बरगडी}} (\emph{bargaḍī}) is a Marathi noun whose origin is not settled.
\end{cousinweb}

\begin{grammarlens}[title={What you've built}]
One more word for describing what you lived and felt.
\end{grammarlens}

\subsection*{Your turn}
\begin{itemize}
\raggedright
  \item \emph{Say it:} \emph{bargaḍī}
  \item \emph{Say it:} \emph{bargaḍī}, once more
  \item \emph{Say it:} \emph{paḍlyāmuḷe tyācī ek bargaḍī moḍlī.} — He broke a rib when he fell.
  \item \emph{Recall:} say the Marathi for the skull, then the Marathi for the spine, then say \emph{bargaḍī} again
\end{itemize}

\begin{tcolorbox}[breakable,colback=teal!4,colframe=teal!35!black,title={Before you move on}]
What does \mr{बरगडी} mean? (\textquotedblleft{}A rib\textquotedblright{}.) Say it once more.
\end{tcolorbox}
\section[\texorpdfstring{\mr{आतडे} (ātḍe)}{आतडे (ātḍe)}]{\mr{आतडे} (ātḍe) — the gut; an intestine}

\label{lesson:MR-C478-atde}



\begin{quote}
Before the new one: say the Marathi for the spine, then the Marathi for a rib.
\end{quote}

\subsection*{You'll want to know: \mr{आतडे}}
\textbf{\mr{आतडे}} — \emph{ātḍe} — \textquotedblleft{}the gut; an intestine\textquotedblright{}.

\begin{quote}
\mr{पोटात} \mr{आतडी} \mr{असतात}.
\end{quote}

\begin{quote}
\emph{poṭāt ātḍī astāt.}
\end{quote}

\begin{quote}
The intestines are in the belly.
\end{quote}

Neuter. Plural: \mr{आतडी}, \emph{ātḍī}.

Neuter; the plural \mr{आतडी} is the intestines.

\begin{cousinweb}
\textbf{\mr{आतडे}} (\emph{ātḍe}) is a tadbhava, inherited from Sanskrit \emph{āntra}, \textquotedblleft{}entrail\textquotedblright{}; CDIAL lists Marathi \emph{ā̃tḍẽ}, \textquotedblleft{}entrail\textquotedblright{} (CDIAL 1182).
\end{cousinweb}

\begin{grammarlens}[title={What you've built}]
One more word for describing what you lived and felt.
\end{grammarlens}

\subsection*{Your turn}
\begin{itemize}
\raggedright
  \item \emph{Say it:} \emph{ātḍe}
  \item \emph{Say it:} \emph{ātḍe}, once more
  \item \emph{Say it:} \emph{poṭāt ātḍī astāt.} — The intestines are in the belly.
  \item \emph{Recall:} say the Marathi for the spine, then the Marathi for a rib, then say \emph{ātḍe} again
\end{itemize}

\begin{tcolorbox}[breakable,colback=teal!4,colframe=teal!35!black,title={Before you move on}]
What does \mr{आतडे} mean? (\textquotedblleft{}The gut; an intestine\textquotedblright{}.) Say it once more.
\end{tcolorbox}
\section[\texorpdfstring{\mr{सांधा} (sāndhā)}{सांधा (sāndhā)}]{\mr{सांधा} (sāndhā) — a joint}

\label{lesson:MR-C478-sandha}



\begin{quote}
Before the new one: say the Marathi for a rib, then the Marathi for the gut.
\end{quote}

\subsection*{You'll want to know: \mr{सांधा}}
\textbf{\mr{सांधा}} — \emph{sāndhā} — \textquotedblleft{}a joint\textquotedblright{}.

\begin{quote}
\mr{गुडघा} \mr{हा} \mr{मोठा} \mr{सांधा} \mr{आहे}.
\end{quote}

\begin{quote}
\emph{guḍghā hā moṭhā sāndhā āhe.}
\end{quote}

\begin{quote}
The knee is a big joint.
\end{quote}

Masculine. Plural: \mr{सांधे}, \emph{sāndhe}.

Masculine; also a seam or a join.

\begin{cousinweb}
\textbf{\mr{सांधा}} (\emph{sāndhā}) is a tadbhava: CDIAL lists Marathi \emph{sā̃dhā}, \textquotedblleft{}joint, knuckle, knot\textquotedblright{}, under Sanskrit \emph{saṁdhi}, \textquotedblleft{}joining\textquotedblright{} (CDIAL 12913).
\end{cousinweb}

\begin{grammarlens}[title={What you've built}]
One more word for describing what you lived and felt.
\end{grammarlens}

\subsection*{Your turn}
\begin{itemize}
\raggedright
  \item \emph{Say it:} \emph{sāndhā}
  \item \emph{Say it:} \emph{sāndhā}, once more
  \item \emph{Say it:} \emph{guḍghā hā moṭhā sāndhā āhe.} — The knee is a big joint.
  \item \emph{Recall:} say the Marathi for a rib, then the Marathi for the gut, then say \emph{sāndhā} again
\end{itemize}

\begin{tcolorbox}[breakable,colback=teal!4,colframe=teal!35!black,title={Before you move on}]
What does \mr{सांधा} mean? (\textquotedblleft{}A joint\textquotedblright{}.) Say it once more.
\end{tcolorbox}
